\documentclass[10pt,a4paper]{amsart}
\usepackage[margin=19mm]{geometry}
\usepackage{amsmath,amssymb,mathtools}
\usepackage[scr=rsfs,cal=euler]{mathalfa}
\usepackage{xfrac}
\usepackage[unicode,hidelinks,bookmarks=false]{hyperref}
\newcommand{\ie}{i.e.\ }
\newcommand{\cf}{cf.\ }
\newcommand{\resp}{respectively\ }
\newcommand{\Subsection}{\subsection}
\providecommand{\nobreakdash}{\nobreak\hbox{-}}
\setlength{\emergencystretch}{2em}
\multlinegap0pt
\usepackage{amssymb}
\usepackage{bm}
\newtheorem{lemma}{Lemma}
\title{Propagation dynamics of an acid-mediated invasion model with degenerate tumor diffusion}
\author{Xinyue Cao, Quentin Griette, Xiong Li, Yang Wang}
\date{Source: arXiv:2607.13944v2 (CC BY 4.0)}
\begin{document}
\maketitle

\noindent\emph{Source and licence.} Propagation dynamics of an acid-mediated invasion model with degenerate tumor diffusion. Version arXiv:2607.13944v2 (CC BY 4.0), distributed by arXiv under the Creative Commons Attribution 4.0 International licence, http://creativecommons.org/licenses/by/4.0/, as recorded in the arXiv licence metadata for this exact version and shown on the abstract page https://arxiv.org/abs/2607.13944v2. Full licence notice: \texttt{licenses/arxiv-2607.13944-CC-BY-4.0.txt}.\par\smallskip

\noindent\textit{Editorial scope.} Counted unit: the article's lemma lem:U-basic (source0.tex lines 490--512). The article's complete original proof of it is delivered with it (source0.tex lines 514--647, one \texttt{proof} environment of 134 lines). No mathematics is written, repaired or re-derived: the statement and the proof are the article's own lines, in the article's own notation, and no retained line is substituted or re-flowed.  Retained uncounted as the source-local context the proof needs: lines 485--487.  Nothing was omitted from any retained span, so the package carries no omission record.  No retained line contains a citation, so the wrapper carries no bibliography and no bibliography entry.  No retained line contains a figure environment, an image-inclusion command or drawing code, so no image is delivered and the package declares no asset dependency.  Editorial formatting only, in addition: an AMS \texttt{amsart} preamble that restates the article's own declarations and macros used by the retained lines, namely the environments \texttt{lemma} on separate counters, together with the standard packages those lines need.  The retained environments print in this document as Lemma 1 in order of appearance, whereas the article prints the same items with the numbering of its own full document.  The complete native source of the article as distributed in the arXiv e-print for this exact version is bundled unchanged as \texttt{sources/205246/source0.tex}.
\begin{equation}\label{eq:T1-def}
T_1(v):=W_v,\qquad v\in\mathcal K_L.
\end{equation}

\begin{lemma}\label{lem:U-basic}
Assume that $d_1>1$. For each $v\in\mathcal K_L$, set
$W_v=T_1(v)$ as in \eqref{eq:T1-def}. Then the equation
\begin{equation}\label{eq:Uv-problem}
\theta U_v'+U_v(1-U_v-d_1W_v)=0,\qquad z\in\mathbb R,
\end{equation}
subject to the boundary condition
\begin{equation}\label{eq:Uv-plus-bc}
U_v(+\infty)=1
\end{equation}
admits a unique solution $U_v\in C^1(\mathbb R)$. This solution satisfies
\begin{equation*}
0<U_v(z)<1,\ z\in\mathbb R,\qquad U_v(-\infty)=0,
\end{equation*}
and is given by
\begin{equation}\label{eq:Uv-formula}
U_v(z)=\theta\left(\int_z^{+\infty}e^{-\frac{1}{\theta}\int_z^s(1-d_1W_v(\tau))\,d\tau}\,ds\right)^{-1},\qquad z\in\mathbb R.
\end{equation}
Moreover, for every compact interval $I\subset\mathbb R$, there exists a constant $C_{IU}>0$, independent of $v$, such that
\begin{equation}\label{eq:U-C1}
\|U_v\|_{C^1(I)}\le C_{IU}.
\end{equation}
\end{lemma}

\begin{proof}
Since $W_v(+\infty)=0$, we have $1-d_1W_v(z)\to1$ as $z\to+\infty$, and hence there exists $R_1>0$ such that
\begin{equation*}
1-d_1W_v(y)\ge\frac12,\qquad y\ge R_1.
\end{equation*}
For each fixed $z\in\mathbb R$, set $R_z:=\max\{z,R_1\}$. Then, for $s\ge R_z$,
\begin{align*}
\int_z^s(1-d_1W_v(\tau))\,d\tau
&=\int_z^{R_z}(1-d_1W_v(\tau))\,d\tau+
\int_{R_z}^s(1-d_1W_v(\tau))\,d\tau\\
&\ge \int_z^{R_z}(1-d_1W_v(\tau))\,d\tau
+\frac12(s-R_z).
\end{align*}
Denote
\begin{equation*}
C_z:=\int_z^{R_z}(1-d_1W_v(\tau))\,d\tau.
\end{equation*}
Since $W_v$ is continuous and $[z,R_z]$ is a finite interval, $C_z$ is finite, and therefore, for $s\ge R_z$,
\begin{equation*}
e^{-\frac{1}{\theta}\int_z^s(1-d_1W_v(\tau))\,d\tau}\le
e^{-\frac{C_z}{\theta}}e^{-\frac{s-R_z}{2\theta}},
\end{equation*}
it follows that
\begin{align*}
\int_{R_z}^{+\infty}e^{-\frac{1}{\theta}\int_z^s(1-d_1W_v(\tau))\,d\tau}\,ds
&\le e^{-\frac{C_z}{\theta}}\int_{R_z}^{+\infty}e^{-\frac{s-R_z}{2\theta}}\,ds\\
&=2\theta e^{-\frac{C_z}{\theta}}<+\infty.
\end{align*}
Moreover, since the integrand is continuous on the finite interval $[z,R_z]$, we also have
\begin{equation*}
\int_z^{R_z}e^{-\frac{1}{\theta}\int_z^s(1-d_1W_v(\tau))\,d\tau}\,ds<+\infty.
\end{equation*}
Combining the two estimates with the positivity of the integrand gives
\begin{equation*}
0<\int_z^{+\infty}e^{-\frac{1}{\theta}\int_z^s(1-d_1W_v(\tau))\,d\tau}\,ds<+\infty.
\end{equation*}
Thus the right-hand side of \eqref{eq:Uv-formula} is well defined and positive. We define $U_v$ by \eqref{eq:Uv-formula}, or equivalently,
\begin{equation}\label{eq:Uv-reciprocal-formula}
\frac{1}{U_v(z)}=\frac{1}{\theta}\int_z^{+\infty}e^{-\frac{1}{\theta}\int_z^s(1-d_1W_v(\tau))\,d\tau}\,ds,\qquad z\in\mathbb R.
\end{equation}

We now verify that \eqref{eq:Uv-reciprocal-formula} indeed gives a positive solution. Since $W_v\in C^2(\mathbb R)$, differentiating \eqref{eq:Uv-reciprocal-formula} gives
\begin{equation}\label{eq:reciprocal-equation-short}
\left(\frac{1}{U_v}\right)'-\frac{1-d_1W_v}{\theta}\frac{1}{U_v}=-\frac{1}{\theta}.
\end{equation}
Indeed, the above differentiation is justified by the exponential decay of the integrand at $+\infty$, locally uniformly with respect to $z$. Since $U_v>0$, multiplying \eqref{eq:reciprocal-equation-short} by $-\theta U_v^2$ yields
\begin{equation*}
\theta U_v'+U_v(1-U_v-d_1W_v)=0,
\end{equation*}
which is exactly \eqref{eq:Uv-problem}. Hence $U_v\in C^1(\mathbb R)$.

It remains to verify the boundary condition \eqref{eq:Uv-plus-bc}. From \eqref{eq:Uv-reciprocal-formula}, we rewrite
\begin{equation*}
\frac{1}{U_v(z)}=\frac{1}{\theta}e^{\frac{1}{\theta}\int_0^z(1-d_1W_v(\tau))\,d\tau}\int_z^{+\infty}e^{-\frac{1}{\theta}\int_0^s(1-d_1W_v(\tau))\,d\tau}\,ds.
\end{equation*}
Equivalently,
\begin{equation*}
U_v(z)=\theta\frac{e^{-\frac{1}{\theta}\int_0^z(1-d_1W_v(\tau))\,d\tau}}{\int_z^{+\infty}e^{-\frac{1}{\theta}\int_0^s(1-d_1W_v(\tau))\,d\tau}\,ds}.
\end{equation*}
Since $1-d_1W_v(z)\to1$ as $z\to+\infty$, and by the estimate above the denominator is the tail of an integrable function, both the numerator and the denominator tend to zero as $z\to+\infty$. Applying L'Hospital's rule, we obtain
\begin{align*}
\lim_{z\to+\infty}U_v(z)
&=\theta\lim_{z\to+\infty}\frac{e^{-\frac{1}{\theta}\int_0^z(1-d_1W_v(\tau))\,d\tau}}{\int_z^{+\infty}e^{-\frac{1}{\theta}\int_0^s(1-d_1W_v(\tau))\,d\tau}\,ds
}\\
&=\theta\lim_{z\to+\infty}\frac{-\frac{1}{\theta}(1-d_1W_v(z))e^{-\frac{1}{\theta}\int_0^z(1-d_1W_v(\tau))\,d\tau}}{-e^{-\frac{1}{\theta}\int_0^z(1-d_1W_v(\tau))\,d\tau}
}\\
&=\lim_{z\to+\infty}(1-d_1W_v(z))=1.
\end{align*}
Thus $U_v(+\infty)=1$.

Next, since $W_v(z)>0$ for every $z\in\mathbb R$,  we have
\begin{equation*}
\int_z^{+\infty}e^{-\frac{1}{\theta}\int_z^s(1-d_1W_v(\tau))\,d\tau}\,ds>
\int_z^{+\infty}e^{-(s-z)/\theta}\,ds=\theta.
\end{equation*}
By \eqref{eq:Uv-formula}, we obtain
\begin{equation*}
0<U_v(z)<1,\qquad z\in\mathbb R.
\end{equation*}
In particular, by hypothesis $(H_1)$,
\begin{equation*}
D(U_v(z))>0,\qquad z\in\mathbb R.
\end{equation*}

To prove the left-end limit, recall that $W_v(-\infty)=1$ and $d_1>1$. Then there exist constants $R_0<0$ and $\delta_0>0$ such that
\begin{equation*}
1-d_1W_v(z)\le-\delta_0,\qquad z\le R_0.
\end{equation*}
For $z<R_0$, it follows from \eqref{eq:Uv-reciprocal-formula} that
\begin{align*}
\frac{1}{U_v(z)}
&\ge\frac{1}{\theta}\int_z^{R_0}e^{-\frac{1}{\theta}\int_z^s(1-d_1W_v(\tau))\,d\tau}\,ds\\
&\ge\frac{1}{\theta}\int_z^{R_0}e^{\frac{\delta_0}{\theta}(s-z)}\,ds
=\frac{1}{\delta_0}\left(e^{\frac{\delta_0}{\theta}(R_0-z)}-1\right).
\end{align*}
Taking $z\to-\infty$, we get $1/U_v(z)\to+\infty$, and therefore
\begin{equation*}
U_v(-\infty)=0.
\end{equation*}

Finally, from \eqref{eq:Uv-problem}, together with $0<U_v<1$ and $0<W_v<1$, we have
\begin{equation*}
|U_v'(z)|
=\frac{1}{\theta}U_v(z)|U_v(z)+d_1W_v(z)-1|
\le\frac{1}{\theta}U_v(z)\bigl(U_v(z)+d_1W_v(z)+1\bigr)
\le\frac{d_1+2}{\theta},
\end{equation*}
for all $z\in\mathbb R$. Hence, for every compact interval $I\subset\mathbb R$, there exists a constant $C_{IU}>0$, independent of $v$, such that
\begin{equation*}
\|U_v\|_{C^1(I)}\le C_{IU}.
\end{equation*}
This proves \eqref{eq:U-C1}.

It remains to prove uniqueness. Let $\widehat U_v\in C^1(\mathbb R)$ be another solution of \eqref{eq:Uv-problem} satisfying \eqref{eq:Uv-plus-bc}. Since $\widehat U_v(+\infty)=1$, we have $\widehat U_v>0$ for all sufficiently large $z$. Indeed, if $\widehat U_v$ had a finite zero, then the local uniqueness theorem for ordinary differential equations applied to \eqref{eq:Uv-problem} would imply $\widehat U_v\equiv0$, contradicting $\widehat U_v(+\infty)=1$. Hence $\widehat U_v(z)>0$ for all $z\in\mathbb R$.

For such a positive solution $\widehat U_v$, division by $\widehat U_v^2$ gives
\begin{equation*}
\left(\frac{1}{\widehat U_v}\right)'
-\frac{1-d_1W_v}{\theta}\frac{1}{\widehat U_v}
=-\frac{1}{\theta}.
\end{equation*}
Solving this linear equation on $[z,R]$ yields
\begin{equation*}
\frac{1}{\widehat U_v(z)}
=\frac{1}{\widehat U_v(R)}e^{-\frac{1}{\theta}\int_z^R(1-d_1W_v(\tau))\,d\tau}
+\frac{1}{\theta}\int_z^Re^{-\frac{1}{\theta}\int_z^s(1-d_1W_v(\tau))\,d\tau}\,ds.
\end{equation*}
Since $W_v(+\infty)=0$ and $\widehat U_v(+\infty)=1$, the first term on the right-hand side tends to zero as $R\to+\infty$. Letting $R\to+\infty$, we obtain
\begin{equation*}
\frac{1}{\widehat U_v(z)}=\frac{1}{\theta}\int_z^{+\infty}e^{-\frac{1}{\theta}\int_z^s(1-d_1W_v(\tau))\,d\tau}\,ds.
\end{equation*}
This coincides with \eqref{eq:Uv-reciprocal-formula}; hence
$1/\widehat U_v=1/U_v$ on $\mathbb R$. Since both functions are positive, we conclude that $\widehat U_v\equiv U_v$, and uniqueness follows.
\end{proof}

\end{document}
